% !TeX root = ../main.tex
\begin{frame}{01 研究背景——支架失效机制}
  \begin{SlideCanvas}
    \BodyAt{55}{84}{850}{支架内再狭窄（ISR）指植入后管腔再次变窄。机械状态与组织反应共同影响支架的长期功能。}
    \PicAt{55}{158}{490}{285}{assets/scai_restenosis.pdf}
    \CaptionAt{55}{452}{485}{冠脉造影、OCT与相应机制示意}
    \GroupHeadingAt{565}{153}{335}{支架失效因素}
    \draw[BICNavy,line width=.8pt] (565,190)--(900,190);
    \BulletPointAt{585}{211}{315}{机械因素}
    \BodyAt{585}{254}{315}{扩张不足、回弹、纵向变形与断裂。}
    \BulletPointAt{585}{327}{315}{生物学因素}
    \BodyAt{585}{370}{315}{内膜增生、新生动脉粥样硬化等。}
    \CaptionAt{585}{432}{315}{ISR：in-stent restenosis\\OCT：光学相干断层成像}
    \Sources{\RefSCAI，图1及机制讨论。左图前两列为扩张不足与断裂，后两列为组织性机制。}
  \end{SlideCanvas}
  \BICConclusionBand{扩张状态与结构完整性是\BICBandEmph{可由器件力学设计调节的环节}}
  \note{“再狭窄”和“支架内再狭窄”均为规范医学术语，本页给出便于医工交叉听众理解的解释，不替代按造影或生理指标确定的诊断标准。SCAI共识将扩张不足列为ISR主要机械原因。机械因素、生物反应和治疗条件相互作用，结构优化的医学收益与研究范围集中在挑战页讨论。本页不把长度变化列为所有支架失效的首要原因。}
\end{frame}
